\chapter*{Preface}
\label{chap:Introduction}

Model categories, introduced by Daniel Quillen \cite{Quillen1}, are a remarkably succesful framework for expressing homotopy theoretic ideas in an axiomatic way. A Quillen model category (${\cat M}$, ${\cat W}$, ${\cat Cof}, {\cat Fib}$) consists of a category ${\cat M}$, and three distinguished classes of maps: the weak equivalences ${\cat W}$, the cofibrations ${\cat Cof}$ and the fibrations ${\cat Fib}$, subject to a list of axioms (\definitionref{defn:modelcat}). 

Let us fix some notation. If ${\cat D}$ is a small category we denote ${\cat M}^{\cat D}$ the category of ${\cat D}$-diagrams in ${\cat M}$. For a functor $u \colon {\cat D}_1 \ra {\cat D}_2$ of small categories we denote $u^* \colon {\cat M}^{{\cat D}_2} \ra {\cat M}^{{\cat D}_1}$ the functor defined by $(u^*X)_{d_1} = X_{ud_1}$ for objects $d_1 \eps {\cat D}_1$.  For a category with a class of weak equivalences $({\cat M}, {\cat W})$ we denote ${\bf ho}{\cat M} = {\cat M}[{\cat W}^{-1}]$ its homotopy category, and we will always consider the weak equivalences on ${\cat M}^{\cat D}$ to be {\it pointwise}, i.e. $f$ is a weak equivalence in ${\cat M}^{\cat D}$ if $f_d$ is a weak equivalence in ${\cat M}$ for all objects $d \eps {\cat D}$.

A Quillen model category ${\cat M}$ admits homotopy pushouts and homotopy pullbacks. These are the total left (resp. right) derived functors of the pushout (resp. pullback) functor in ${\cat M}$. If ${\cat M}$ is pointed one can construct the homotopy cofiber and fiber of a map in ${\bf ho}{\cat M}$, and the suspension and loop space of an object in ${\bf ho}{\cat M}$. One can then form the cofibration sequence and the fibration sequence of a map in ${\bf ho}{\cat M}$.

Furthermore, a Quillen model category ${\cat M}$ admits all small homotopy colimits and limits (also called the homotopy left and right Kan extensions). For a functor $u \colon {\cat D}_1 \ra {\cat D}_2$ of small categories, the homotopy colimit ${\bf L}\colim^u \colon {\bf ho}({\cat M}^{{\cat D}_1}) \ra {\bf ho}({\cat M}^{{\cat D}_2})$ is the total left derived functor of the colimit functor $\colim^u$, and is left adjoint to ${\bf ho}u^* \colon {\bf ho}({\cat M}^{{\cat D}_2}) \ra {\bf ho}({\cat M}^{{\cat D}_1})$. Dually, the homotopy limit ${\bf R}\lim^u$ is the total right derived functor of $\lim^u$, and is right adjoint to ${\bf ho}u^*$.

For the general construction of homotopy colimits, we recommend the work of Dwyer, Kan, Hirschhorn and Smith \cite{Dwyer-Kan-Hirschhorn-Smith}, Hirschhorn \cite{Hirschhorn}, Weibel \cite{Weibel} (after Thomason's unpublished notes), Chach\'{o}lski and Scherer \cite{ChacScherer} and Cisinski \cite{Cisinski4}. Each of these references presents a different perspective on homotopy (co)limits. It is apparent from these sources that in a Quillen model category the homotopy colimits can actually be constructed just manipulating weak equivalences and cofibrations. 

\bigskip

To summarize, homotopy colimits are the total left derived functors of the colimit, so their {\it definition} requires just the presence of weak equivalences. The {\it existence} of homotopy colimits can be proved in presence of the Quillen model category axioms, where their {\it construction} can be performed using just cofibrations and weak equivalences.

It is therefore natural to ask if one can simplify Quillen's axioms, and separate a minimal set of axioms required by cofibrations and weak equivalences in order to still be able to prove existence of homotopy colimits. Along the way, this leads us to rethink the role of cofibrations in abstract homotopy theory.

\bigskip

In a very influential paper, Ken Brown \cite{Brown} formalized some of these obsevations by defining categories of fibrant objects and working out in detail their properties. Reversing arrows, one defines categories of cofibrant objects, and Brown's work carries over by duality to categories of cofibrant objects. 

A category of cofibrant objects (${\cat M}$, ${\cat W}$, ${\cat Cof}$) consists of a category ${\cat M}$, the class of weak equivalences ${\cat W}$ and the class of cofibrations ${\cat Cof}$, subject to a list of axioms (the duals of the axioms of \cite{Brown}). These axioms require in particular all objects to be cofibrant.

For a pointed category of cofibrant objects, Brown was able to construct homotopy cofibers of maps, suspensions of objects and the cofibration sequence of a map. These constructions exist in dual form for categories of fibrant objects.

\bigskip

Building on Brown's work, Don Anderson \cite{Anderson1} extended Brown's axioms for a category of cofibrant objects by dropping the requirement that all objects be cofibrant. Anderson called the categories defined by his new axioms {\it left homotopical}; our text changes terminology and calls them {\it Anderson-Brown-Cisinski cofibration categories} (or just cofibration categories for simplicity). The cofibration category axioms we use are slightly more general than Anderson's.

Anderson's main observation was that the cofibration category axioms on ${\cat M}$ suffice for the construction of a left adjoint of ${\bf ho}u^*$, for any functor $u$ of small categories. It is implicit in his work that the left adjoint of ${\bf ho}u^*$ is a left derived of $\colim^u$.

Unfortunately for the history of this subject, Anderson's paper \cite{Anderson1} contains statements but omits proofs, and has a title (``Fibrations and Geometric Realizations``) that does not reflect the generality of his work. Also, Anderson quit mathematics shortly after his paper was published, the proofs of \cite{Anderson1} got lost and as a result his whole theory laid dormant for twenty five years.

\bigskip

We can be grateful to Denis-Charles Cisinski \cite{Cisinski2}, \cite{Cisinski4} for bringing back to light Brown and Anderson's ideas. Cisinski simplifies Anderson's arguments, and provides for cofibration categories a complete construction of homotopy colimits along functors $u \colon {\cat D}_1 \ra {\cat D}_2$ with ${\cat D}_1$ finite and direct. 

Cisinski has also worked out the construction of homotopy colimits along arbitrary functors $u$ of small categories, as well as the end result regarding the derivability of cofibration categories (our \chapterref{chap:derivators}). While this part of his work remains unpublished, he was kind enough to share with me its outline. I would like to thank him for suggesting the correct formulation of axioms CF5-CF6, and for patiently explaining to me the finer points of excision. 

\bigskip

The goal of these notes is then to work out an account of homotopy colimits from the axioms of an Anderson-Brown-Cisinski cofibration category, and show that they satisfy the axioms of a left Heller derivator. There are a number of properties of homotopy colimits that are a formal consequence of the left Heller derivator axioms, but they are outside of the scope of our text. We will instead try to investigate the relation with the better-known Quillen model categories, and compare with other axiomatizations that have been proposed for cofibration categories. 

While some of the proofs we propose may be new, the credit for this theory should go entirely to Brown, Anderson and Cisinski. It was our choice in this text to make use of approximation functors and abstract Quillen equivalences, and for that we were influenced by the work of Dwyer, Kan, Hirschhorn and Smith \cite{Dwyer-Kan-Hirschhorn-Smith}. Our treatment of direct categories bears the influence of Daniel Kan's theory of Reedy categories outlined in \cite{Dwyer-Kan-Hirschhorn}, \cite{Hovey} and \cite{Hirschhorn}.

\bigskip

I would like to thank Haynes Miller and Daniel Kan, my mentors in abstract homotopy, for their gracious support and encouragement. I am grateful to Denis-Charles Cisinski, Philip Hirschhorn and Haynes Miller for the conversations we had on the subject of this text.

\subsection*{Comparing ABC and Quillen model categories}
\label{subsec:compareabcquillen}

Anderson-Brown-Cisinski (ABC) cofibration categories (${\cat M}$, ${\cat W}$, ${\cat Cof}$) are defined in \sectionref{sec:CofibrationCategoriesFibrationCategoriesAxioms}. We will make a distinction between ABC precofibration categories (satisfying axioms CF1-CF4) and ABC cofibration categories (satisfying the full set of axioms CF1-CF6). 

ABC fibration categories (${\cat M}$, ${\cat W}$, ${\cat Fib}$) are defined by duality, and ABC model categories (${\cat M}$, ${\cat W}$, ${\cat Cof}$, ${\cat Fib}$) by definition carry an ABC cofibration and fibration category structure.

Any Quillen model category is an ABC model category, but the class of ABC cofibration categories behaves differently in some respects:
\begin{enumerate}
\item
If (${\cat M}$, ${\cat W}$) admits a Quillen model category structure, the choice of cofibrations determines the fibrations, and viceversa. This is not true of an ABC cofibration category structure, where the choice of cofibrations is completely independent of the choice of fibrations.
\item
If (${\cat M}$, ${\cat W}$, ${\cat Cof}$) is an ABC cofibration category, then so is (${\cat M}^{\cat D}$, ${\cat W}^{\cat D}$, ${\cat Cof}^{\cat D}$) for a small category ${\cat D}$. Quillen model categories have this property only in particular cases, for example if they are cofibrantly generated.
\item
If ${\cat M}$ is a locally small Quillen model category, then ${\bf ho}{\cat M}$ is also locally small. ABC cofibration categories do not have this property.
\item
If (${\cat M}$, ${\cat W}$, ${\cat Cof}$, ${\cat Fib}$) is a Quillen model category, then ${\cat W}$ is saturated ${\cat W} = \overline{\cat W}$. Similarly, if (${\cat M}$, ${\cat W}$, ${\cat Cof}$) is an ABC cofibration category then ${\cat W}$ is saturated. If (${\cat M}$, ${\cat W}$, ${\cat Cof}$) is only a precofibration category, then ${\cat W}$ is not necessarily saturated, but (${\cat M}$, $\overline{\cat W}$, ${\cat Cof}$) is still an ABC precofibration category.
\item
If (${\cat M}$, ${\cat W}$, ${\cat Cof}$) is a left proper ABC cofibration category, one has a {\it maximal} left proper CF1-CF4 cofibration category structure on (${\cat M}$, ${\cat W}$), taking as cofibrations the left proper maps ${\cat PrCof}$. In this context, the weak equivalences therefore determine the left proper maps as a preferred class of CF1-CF4 cofibrations.
\end{enumerate}

\subsection*{Outline of the text}
\label{subsec:outline}

We start \chapterref{chap:CofibrationAndFibrationCategories} with the ABC (pre)cofibration category axioms, we proceed with their elementary properties, and end the chapter with a discussion of proper ABC cofibration categories. For a cofibration category ${\cat M}$, we will denote ${\cat M}_{cof}$ its full subcategory of cofibrant objects.

Going to \chapterref{chap:modcatcatcofobj}, we compare ABC cofibration categories with other axiomatic systems that have been proposed for categories with cofibrations. Most notably, we show that any Quillen model category is an ABC model category.

Topological spaces, with homotopy equivalences as weak equivalences, Hurewicz cofibrations as cofibrations, and Hurewicz fibrations as fibrations, form an ABC model category. Complexes of objects in an abelian category ${\cat A}$, with quasi-isomorphisms as weak equivaleces, and monics as cofibrations form an ABC precofibration category, which satisfies CF5 if ${\cat A}$ satisfies the Grothendieck axiom AB4, and satisfies CF6 if ${\cat A}$ satisfies the Grothendieck axiom AB5. These examples are explained in \chapterref{chap:examples}.

In \chapterref{chap:kanextensions} we recall the language of 2-categories and that of Kan extensions. This is a prerequisite for \chapterref{chap:catweq}, where we introduce axiomatically the {\it left approximation} functors $t \colon ({\cat M}^{'}, {\cat W}^{'}) \ra ({\cat M}, {\cat W})$ between two category pairs. The Approximation \theoremref{thm:approxthm} states that left approximation functors induce an equivalence of categories ${\bf ho}t \colon {\bf ho}{\cat M}^{'} \ra {\bf ho}{\cat M}$. We prove an existence result for total derived functors (\theoremref{thm:generalexistenceleftderived}), and a Quillen-type adjunction property between the total derived functors of an adjoint pair of functors (\theoremref{thm:generalexistencetotalderivedadjoint2}). 

In \chapterref{chap:homotcat}, we take as model the inclusion ${\cat M}_{cof} \ra {\cat M}$ and define {\it cofibrant approximation} functors $t \colon {\cat M}^{'} \ra {\cat M}$ between precofibration categories. Cofibrant approximation functors are in particular left approximations, and by the Approximation \theoremref{thm:approxthm} they induce an equivalence of categories ${\bf ho}t$. In particular, this shows that ${\bf ho}{\cat M}_{cof} \ra {\bf ho}{\cat M}$ is an equivalence of categories.

In the second half of \chapterref{chap:homotcat} we then recall the theory of cylinders and left homotopic maps in a precofibration category, and develop the formalism of homotopy calculus of fractions. Going to \chapterref{chap:homcalsappl}, we give some applications of the formalism of homotopy calculus of fractions.

In \chapterref{chap:constrcat}, we recall over and under categories, and elementary properties of limits and colimits.

In \chapterref{chap:limcolimcofinal}, we reach our main objective. We define the homotopy colimit ${\bf L}\colim^u \colon {\bf ho}({\cat M}^{{\cat D}_1})$ $\ra$ ${\bf ho}({\cat M}^{{\cat D}_2})$ as a left derived functor, and show that it exists and it is left adjoint to ${\bf ho}u^* \colon {\bf ho}({\cat M}^{{\cat D}_2})$ $\ra$ ${\bf ho}({\cat M}^{{\cat D}_1})$. We show that if ${\cat D}$ is a small category and (${\cat M}$, ${\cat W}$, ${\cat Cof}$) is a cofibration category, then the diagram category ${\cat M}^{\cat D}$ is again a cofibration category, with pointwise weak equivalences and pointwise cofibrations.

Finally, in \chapterref{chap:derivators} we recall the notion of a {\it left Heller derivator}, and show that the homotopy colimits in a cofibration category satisfy the axioms of a left Heller derivator. The purpose of the last chapter is simply to assert the results of \chapterref{chap:limcolimcofinal} within the axiomatic language of derivators.

